\section{线性近似与极值问题}
\label{sec:02-Calculus/03-Differentiation/05}

朋友发来一条消息："帮我心算 \(\sqrt{4.1}\)，身边没有计算器。"你知道 \(\sqrt{4}=2\)，但多出来的 \(0.1\) 会贡献多少？

最容易想到的招数：既然 \(\sqrt{4}=2\)，\(\sqrt{9}=3\)，在 \(4\) 和 \(9\) 之间线性内插——\(\sqrt{4.1}\) 大约在 \(2\) 到 \(3\) 的五十分之一处，估算 \(2.02\)。但这是拿整段 \([4,9]\) 的平均斜率去套 \([4,4.1]\) 上的变化——太粗糙。

更好的招数：导数告诉你函数在一个点上的瞬时变化速度。\(f(x)=\sqrt{x}\) 在 \(x=4\) 的导数是

\[
f'(4)=\frac{1}{2\sqrt{4}}=\frac14=0.25.
\]

如果这个变化速度在从 \(4\) 到 \(4.1\) 的短短 \(0.1\) 的距离内大致维持不变，增量大约是 \(0.25\times0.1=0.025\)，所以

\[
\sqrt{4.1}\approx 2+0.025=2.025.
\]

用计算器验证：\(\sqrt{4.1}\approx2.024845\ldots\)，差距不到万分之二。这个精度值得追究——不是碰巧。

\subsection{可微的精确含义：剩下的误差比步长缩得更快}

刚才的做法在几何上等于用切线代替曲线。把函数在 \(x=a\) 附近展开：

\[
f(a+h)=f(a)+f'(a)h+r(h),
\]

其中 \(r(h)\) 是真实值与切线估计之差。函数在 \(a\) 可微的条件精确地说，正是这个残差不仅随 \(h\to0\) 趋于零，而且比 \(h\) 缩得还快：

\[
\lim_{h\to0}\frac{r(h)}{h}=0.
\]

这句话拆开来看：当步长减半，线性近似的主要误差（绝对值）大约减为原来的四分之一——因为误差是 \(h\) 的高阶项，主项与 \(h^2\) 成正比。用 \(\sqrt{4.1}\) 的相邻点验证一下这个规律。取 \(h=0.1\)：

\[
\sqrt{4.1}\approx 2.025,\quad\text{实际}\approx 2.02485,\quad \text{误差}\approx -1.55\times 10^{-4}.
\]

取 \(h=0.01\)：\(\sqrt{4.01}\approx 2.0025\)，实际 \(\approx 2.002498\)，误差 \(\approx -2\times 10^{-6}\)。步长十分之一，误差约百分之一——确实比线性关系消失得更快。

由此，线性化

\[
L(x)=f(a)+f'(a)(x-a)
\]

不只是"切线方程"——它是函数在 \(a\) 点附近误差相对于步长可忽略的一阶主导模型。上一节用中值定理得到过严格上界 \(|\sqrt{4.1}-2|\le0.025\)，那是硬约束；线性近似给出的是中心估计 \(2.025\)。两组信息配合使用：点估计加误差边界，形成一个可信区间。

微分记号 \(dy=f'(x)\,dx\) 记录的正是同一件事：输入微小偏移 \(dx\) 经局部线性映射 \(f'(x)\) 缩放后得到的预测输出变化。记号紧凑，但要警惕——不要把它当成两个"无穷小实数"在做随意除法。

\subsection{离展开点一远，线性近似就露出原形}

线性近似是一个局部的好朋友，不是一个全局的可靠代理。试试用 \(x=1\) 处的切线去估计 \(\sqrt{4}\)：\(f(1)=1\)，\(f'(1)=1/2=0.5\)，线性近似给出

\[
\sqrt{4}\approx 1 + 0.5\times(4-1) = 2.5.
\]

真正的 \(\sqrt{4}=2\)。误差 \(0.5\)，惨不忍睹。原因显而易见——\(\sqrt{x}\) 在 \([1,4]\) 上的弯曲弧度完全不能被一条直线捕捉。线性近似只有在你愿意待在展开点附近时才值得信任。多近算近，取决于函数弯曲得有多厉害——二阶导数的大小决定了近似的有效半径。

\subsection{极值从哪来：围栏问题与临界点搜索}

换一个方向。用 100 米围栏靠墙围一个矩形场地，靠墙一侧不设围栏。怎样围面积最大？

先把自然语言翻译成变量和约束。设垂直墙的边长 \(x\)，平行墙的边长 \(y\)。围栏只用三边：

\[
2x+y=100,\qquad x\ge0,\;y\ge0.
\]

目标面积

\[
A(x)=x(100-2x)=100x-2x^2,\qquad 0\le x\le 50.
\]

可以手动试探几个点：\(x=10\) → \(A=800\)；\(x=20\) → \(A=1200\)；\(x=25\) → \(A=1250\)；\(x=30\) → \(A=1200\)；\(x=40\) → \(A=800\)。最大值显然不在两端（面积都归零），而在中间某处。\(25\) 看起来是峰值——但这是描点描出来的，拐个弯问：为什么偏偏是 \(25\)？描点法只能给出候选，不能给出理由。

如果 \(f\) 在一个内部点 \(c\) 取得局部最大值并且可导——从左趋近 \(c\) 时函数在上升，斜率非负；从右趋近时函数在下降，斜率非正。左右斜率必须相等（可导），唯一的出路是 \(f'(c)=0\)。这就是 Fermat 引理：内部极值点若是可导的，导数必为零。

由此得出闭区间 \([a,b]\) 上求全局极值的标准流程：

\begin{enumerate}
\tightlist
\item 找出所有内部临界点——满足 \(f'(x)=0\) 的点，以及导数不存在的点；
\item 把两个端点 \(a,b\) 也加入候选清单；
\item 比较所有候选点上的函数值，最大者为全局最大，最小者为全局最小。
\end{enumerate}

围栏问题：\(A'(x)=100-4x\)，令其为零得 \(x=25\)，此时 \(y=100-2\times25=50\)，面积 \(A=1250\)。端点 \(x=0\) 和 \(x=50\) 面积均为零。无需再猜——\(25\) 就是最优解。

真正困难常常不在求导之后，而在求导之前：选择变量、写出约束、消去多余变量并确定可行区间。如果忘了 \(x\in[0,50]\) 这个物理约束，代数模型会允许负的边长——数学上函数有定义，物理上毫无意义。

换一个约束类型巩固这个流程。用固定面积 \(A=600\) 平方厘米的纸板做一个无盖方盒，四角切掉正方形后折起。切掉的正方形边长 \(x\) 取多少时盒子容积最大？纸板原边长设为 \(a\times b\)，切角后盒子的长宽高分别是 \(a-2x\)、\(b-2x\)、\(x\)。若纸板是正方形 \(a=b=\sqrt{A}\)，容积 \(V(x)=x(\sqrt{A}-2x)^2\)，定义在 \(0\le x\le\sqrt{A}/2\)。求导、找临界、比较端点——流程和围栏问题完全相同，变量选择和域约束同样是建模阶段最需要警觉的步骤。

\subsection{临界点不等于极值点，极值也可能不在内部}

只解 \(f'(x)=0\) 的人迟早摔进两个坑。

第一个坑：\(f(x)=x^3\) 在 \(x=0\)。\(f'(0)=0\)，但函数值在零点左侧为负、右侧为正——这只是一个水平拐点，切线放平但函数照常穿过。导数为零不保证极值，只保证"瞬时停止变化"——停止后可能转向，也可能继续向前。

第二个坑：\(f(x)=|x|\) 在 \(x=0\)。导数不存在，但这里是全局最小值。导数不存在的点同样进入候选清单。

第三个坑更隐蔽：闭区间上的极值可能就落在端点。\(f(x)=x\) 在 \([0,1]\) 上导数恒为 \(1\)，没有内部临界点，最大值却在 \(x=1\)。只解 \(f'=0\) 会漏掉整个候选空间。

一阶导数符号变化提供了一种不依赖二阶导数的判别：穿过临界点时，导数由正变负则是局部极大；由负变正则局部极小；不变号则不是极值。

\subsection{二阶导数告诉你曲率：碗口朝上还是朝下}

一阶导数告诉你函数在上升还是下降。二阶导数告诉你这种上升或下降本身在加速还是减速——图像是弯向上还是弯向下。

\(f''(x)>0\)：斜率递增，曲线呈凹向上（凸函数），切线始终在曲线下方。\(f''(x)<0\)：斜率递减，凹向下，切线在曲线上方。当 \(f'(c)=0\) 且 \(f''(c)>0\)，切线水平且曲线向上弯——局部极小，碗形。\(f''(c)<0\) 推局部极大。

围栏问题中 \(A''(x)=-4<0\) 处处成立，临界点 \(x=25\) 确认为最大值——无需画图，无需描点表，二阶导数一锤定音。

\subsection{二阶导数为零时：回去看一阶符号}

\(f''(c)=0\) 本身给不出结论。比较同一位置的两种可能：

\(f(x)=x^4\) 在 \(x=0\)：\(f'(0)=0\)，\(f''(0)=0\)。但四次方后两侧皆正——这是全局极小点。

\(f(x)=x^3\) 在 \(x=0\)：\(f'(0)=0\)，\(f''(0)=0\)。但三次方在零点变号——只是拐点。

二阶导数都是零，结论截然不同。当 \(f''=0\) 时，退回到一阶导数变号判别，或检查更高阶导数的首个非零项，或回到极值定义本身。

拐点——凹凸性发生改变的点——同样不能只靠解 \(f''(x)=0\) 来锁定。\(f(x)=x^4\) 在 \(x=0\) 处 \(f''(0)=0\)，但两侧都是凹向上，凹凸性从未改变。拐点要求 \(f''\) 穿过零点并确实发生符号反转。

考一个真正拐点的例子：\(f(x)=x^3-3x\)。求二阶导 \(f''(x)=6x\)。在 \(x=0\) 处 \(f''(0)=0\)。检查两侧：\(x<0\) 时 \(f''(x)<0\)（凹向下），\(x>0\) 时 \(f''(x)>0\)（凹向上）。符号的确反转了——\(x=0\) 是一个货真价实的拐点，曲线在此从''下弯''翻成''上弯''。一阶导数 \(f'(x)=3x^2-3=3(x^2-1)\) 在 \(x=0\) 处为 \(-3\neq0\)——拐点可以是切线有坡度的点，不必水平。

\subsection{从一阶近似到 Taylor 展开：不够精就加一阶}

线性近似的精度受限于曲率。需要更贴近曲线的局部形态时，可以继续匹配更高阶的变化信息：

\[
f(a+h)\approx f(a)+f'(a)h+\frac{f''(a)}{2}h^2+\cdots+\frac{f^{(n)}(a)}{n!}h^n.
\]

线性化是 Taylor 展开截断到一阶。二阶项 \(\frac{f''(a)}{2}h^2\) 纠正了线性近似完全忽略的曲率效应：\(\sqrt{4.1}\) 的线性近似 \(2.025\)，加上二阶修正后约为 \(2.02484\)，距离真实值只剩五百万分之一。更高阶项继续收紧。

但 Taylor 多项式用的是展开点 \(a\) 处的导数信息，对远离 \(a\) 的区域不自动保证精度。后续级数单元会系统处理余项的控制和函数是否等于其 Taylor 级数的问题。此刻只需知道：线性近似是方便的第一刀；精度不够时，下一阶在哪里加、代价是什么，方向是清楚的。

在实际使用中，线性近似还有一个经常被忽略的优势：它给你一个容易心算或纸笔验证的边界判断。工程中测到一个量 \(x=4.00\pm0.05\) 并要求计算 \(\sqrt{x}\) 时，用线性近似立得估计区间 \([1.9875,2.0125]\)——不需要函数表也不需要编程。多数现场决策只需要这个一阶精度；二阶以上是留给需要签字的设计评审的。

\subsection{线性近似与极值：同一套导数工具的两种用法}

回顾这一节和上一节走过的路。中值定理用导数串起了区间两端的全局变化——它告诉你''变化量大致等于导数乘步长，而且存在一点精确成立''。线性近似用导数构造了局部的切线模型——它告诉你''在展开点附近，函数的行为由一阶项主导''。极值判别用导数找出了函数停止变化并可能转向的位置——临界点和二阶曲率合作，区分山顶、谷底和平台。

这三件事表面看是三个独立的话题：误差界、函数逼近、最优化。但它们的共用引擎是同一个：导数把一个复杂的非线性函数在局部掰直成线性，然后用这条直线的信息做判决。中值定理保证掰直的过程有严格误差控制；线性近似提供掰直后的预测值；极值分析利用掰直后斜率为零这个特征信号锁定候选点。

导数的威力不在于它能算切线斜率——这个高中就知道了。它的威力在于它把''局部线性''这个简单得荒唐的想法，推到了一个能严格支撑全局判断的高度。

线性近似还在另一个方向上有直接的延伸：如果你不是想逼近函数值，而是想找函数等于零的点——解方程 \(f(x)=0\)。在猜值 \(x_0\) 处做线性近似：\(f(x)\approx f(x_0)+f'(x_0)(x-x_0)\)。令近似为零解出 \(x\)，得到改进的猜测 \(x_1=x_0-f(x_0)/f'(x_0)\)。再以 \(x_1\) 为起点重复——这就是 Newton 法。它每次迭代只做一次线性近似和一次除法，但收敛速度极快。线性近似的思想在这里从''估计一个数''扩展成了''迭代逼近一个解''——函数的局部线性化不只是静态的快照，也可以成为动态搜索的引擎。
